\chapter{What this kit refuses to claim}
\label{app:refused}

Every figure below circulates widely in preparation material for this role.
None of them is in this book, and this appendix names each one and says why.

\judgement{This is the most useful page in the book, and the reason is not
pedantry. A candidate who prepares against a fabricated loop prepares for
nothing; a candidate who quotes an unsourced premium in a salary conversation
is quoting a blog post at somebody who knows their own bands. The cost of these
numbers is not that they are wrong \dash{} several may well be roughly right
\dash{} it is that nobody can tell, and acting on a number you cannot trace is
the exact failure this book spends thirty-eight chapters arguing against.}

\section{The refused list}

\subsection*{``AI engineering roles command a 25--40 per cent premium''}

\textbf{Traced to:} a single post, no methodology, no corpus, no date range, no
region.

\textbf{Why it is refused:} a premium is a comparison between two populations
and nobody defined either. Against which roles, in which market, at which
level, when? A number with none of those attached is not a measurement.

\textbf{What to do instead:} Chapter~\ref{ch:company}. Establish the band's
composition and level span for the specific company, from what they publish.

\subsection*{``Around 70 per cent of senior loops contain no algorithmic coding''}

\textbf{Traced to:} one engineer's account of her own interviews, at a
self-selected set of companies, in one stretch of time.

\textbf{Why it is refused:} it is widely quoted as a base rate and it is one
observer's sample. The account itself is a useful first-hand source and is
cited in this book for something it can support \dash{} the segmentation in
Chapter~\ref{ch:segment}, which is a way of organising a decision rather than a
frequency claim.

\textbf{What to do instead:} Chapter~\ref{ch:loop}. Ask what the stages are.
The answer for the company in front of you is obtainable in one message.

\subsection*{``Fewer than ten per cent of engineers reach Staff''}

\textbf{Traced to:} recruiting content only.

\textbf{Why it is refused:} no population, no ladder, no counting method. Ten
per cent of whom \dash{} engineers at companies with a published ladder, or
everybody? Counted at a point in time or over a career?

\textbf{What to do instead:} Chapter~\ref{ch:staff}. Read two or three
published ladders and find out what the level means where you are applying.

\subsection*{Named companies' interview loops, stage by stage}

\textbf{Traced to:} preparation vendors citing unnamed candidate reports.

\textbf{Why it is refused:} every version found was a vendor page citing
anonymous sources, and some of it is plausibly invented. The danger is specific
and severe: preparing against a fabricated loop is worse than not preparing,
because it produces confidence.

\textbf{What to do instead:} ask. Chapter~\ref{ch:loop} has the four questions
and they are routinely answered.

\subsection*{``You can become a competent AI engineer in two to three months''}

\textbf{Traced to:} one source, restating its own framing.

\textbf{Why it is refused:} competent at what, from what starting point, judged
by whom. It is a plausible figure and it is an assertion.

\textbf{What to do instead:} Chapter~\ref{ch:gap-map} orders the work. How long
it takes depends on where you are starting, which the figure does not ask.

\subsection*{Individual striking outcomes}

\textbf{The shape:} an engineer publishes an artefact and has an offer within
days; a candidate demonstrates a large cost reduction and is hired immediately.

\textbf{Why it is refused:} single anecdotes, selected for being striking, with
no denominator. The interesting question is how many people published a similar
artefact and heard nothing, and that number is not knowable.

\textbf{What survives:} the general practice is reported in a corpus with a
stated method and is used in Chapter~\ref{ch:portfolio}. The dramatic timelines
are not.

\section{Two claims that are in the book, marked}

Two single-source claims survived, because they change what a reader would
prepare and refusing them would cost more than marking them. Both say
\emph{one source} in the sentence itself.

The first is that at more senior levels, one large employer's assisted coding
round is reported to replace the unassisted one rather than join it
(Chapter~\ref{ch:formats}). Two preparation sites carry it and both trace to
the same internal note, so it is one source wearing two hats.

The second is the three-segment split of the market
(Chapter~\ref{ch:segment}), from one engineer's account. It is used as a way to
organise a decision rather than as a fact about the market.

\judgement{Marking a claim is not a way of keeping it while avoiding
responsibility for it. The test applied was whether a reader who discovered it
was wrong would have wasted effort or made a bad decision. For these two the
answer is that they would have asked a company one extra question, which is
cheap.}

\section{The test, for the next one}

\begin{kitmethod}
\textbf{1.} Who is the origin \dash{} a practitioner describing their own work,
a company describing its own process, or a corpus with a stated method? Or is
it a site whose business is preparation material?

\textbf{2.} Does the claim name its population and its method? A premium needs
two populations; a proportion needs a denominator; a loop needs a source.

\textbf{3.} Do the several sources repeating it trace to one origin? This is
the commonest failure and it produces the appearance of consensus from one
post.

\textbf{4.} If it turned out to be wrong, what would you have done differently?
If the answer is ``nothing'', mark it and move on. If the answer is ``spent a
month'', refuse it.
\end{kitmethod}

Chapter~\ref{ch:noise} applies the same test to topics rather than to figures.
It is the same habit, and it is the one this book is most trying to transmit:
\judgement{a claim you cannot trace is not a small problem with a number, it is
a decision you are making on somebody else's unexamined authority.}
